\documentclass[11pt]{article}
\usepackage[margin=1in]{geometry}
\usepackage{times}
\usepackage{hyperref}
\usepackage{booktabs}
\usepackage{graphicx}
\usepackage{amsmath}
\usepackage{natbib}

\title{\textbf{The Conformity Problem: Social Influence and Belief Revision in Multi-Agent LLM Systems}}
\author{[Your Name] \\ [Your Institution / Independent Researcher] \\ \texttt{[your email]}}
\date{}

\begin{document}
\maketitle

% ============================================================
\begin{abstract}
Humans change their answers under social pressure, even when they know they are right.
Do large language models (LLMs) do the same? This paper presents the first systematic
study of conformity behaviour in LLMs, directly inspired by Asch's classic social
conformity experiments. We design a social influence protocol in which a model that
has correctly answered a question is shown a fabricated response from a ``peer agent''
giving the wrong answer at varying levels of confidence and authority. We measure how
often models abandon correct answers in favour of the wrong one across four frontier
LLMs (GPT-4o, Claude Sonnet, Llama~3~70B, and Gemini~1.5~Flash), seven social
influence conditions, and 200 questions from two reasoning benchmarks (GSM8K and
ARC Challenge). We find conformity rates are substantial — models revise correct
answers to wrong ones at rates exceeding 30\% under high-confidence authoritative
wrong agents. We introduce the \textit{social robustness score} as a new evaluation
metric for multi-agent LLM systems, and show that models differ dramatically in their
resistance to peer pressure. Our findings have direct implications for the design of
multi-agent debate frameworks, ensemble methods, and any system where LLMs are
exposed to each other's outputs.
\end{abstract}

% ============================================================
\section{Introduction}
\label{sec:intro}

In 1951, Solomon Asch demonstrated one of social psychology's most unsettling findings:
ordinary people will give obviously wrong answers to simple perceptual questions when
surrounded by confederates who unanimously give the wrong answer \citep{asch1951effects}.
The Asch conformity experiments revealed that social pressure could override clear
perceptual evidence, with roughly one third of responses conforming to the incorrect
majority view.

As large language models (LLMs) are increasingly deployed in multi-agent configurations
--- debate frameworks \citep{du2023improving}, collaborative coding assistants,
ensemble reasoners, and agentic pipelines where models critique each other's outputs
--- a natural and largely unasked question arises: do LLMs exhibit analogous
conformity behaviour? If a model correctly answers a question, then is exposed to a
confident but wrong response from a peer agent, does it revise its correct answer?

This question is not merely theoretical. Multi-agent LLM systems are premised on the
idea that inter-agent interaction improves output quality. But if agents conform to
each other's errors, the system may silently propagate and amplify mistakes rather
than correct them. A confidently wrong agent could corrupt an entire ensemble,
precisely because its confidence signals authority to other agents trained to be
responsive to human and AI feedback.

We present the first controlled study of LLM conformity behaviour. Our contributions
are as follows:

\begin{itemize}
    \item \textbf{A social influence protocol for LLMs}, directly adapted from the
    Asch paradigm, that measures conformity under varying levels of peer confidence,
    authority framing, and majority pressure.

    \item \textbf{The social robustness score}, a new metric capturing a model's
    resistance to peer pressure across conditions, enabling the first ranking of
    frontier LLMs by social robustness.

    \item \textbf{Empirical evidence of substantial conformity} in all tested models,
    with conformity rates exceeding 30\% under high-confidence authoritative wrong
    agents — and significant variation across models, conditions, and question types.

    \item \textbf{A characterisation of the authority premium} — the additional
    conformity induced by expert or advanced-model framing beyond a neutral peer,
    quantifying how authority signals modulate LLM belief revision.
\end{itemize}

% ============================================================
\section{Related Work}
\label{sec:related}

\subsection{Sycophancy in LLMs}
The most closely related line of work studies sycophancy — the tendency of LLMs to
agree with user assertions even when incorrect \citep{perez2022sycophancy,
sharma2023towards}. Sycophancy research typically examines human-to-model influence:
a user pushes back on a model's answer, and the model capitulates. Our work is
distinct in two ways: we study \textit{agent-to-agent} influence (not human-to-model),
and we use a controlled protocol where the ground truth is known, enabling precise
measurement of when conformity leads to error.

\subsection{Multi-Agent LLM Systems}
Recent work has explored multi-agent debate as a mechanism for improving factuality
and reasoning \citep{du2023improving, liang2023encouraging}. These papers generally
find that debate improves aggregate accuracy, but do not study failure modes where
agents converge on wrong answers. \citet{xiong2023llms} examine uncertainty
expression in LLMs but do not study how expressed confidence from one agent affects
another. Our work fills this gap.

\subsection{Prompt Sensitivity and Robustness}
A growing body of work documents that LLM outputs are sensitive to surface-level
prompt variations \citep{mizrahi2023state, chen2025calibration}. This work
establishes that models are not robust to input perturbations, but focuses on
single-agent settings. Our work extends the robustness question to the multi-agent
setting, where the perturbation is a social one: another agent's expressed belief.

\subsection{Social Psychology and AI}
The application of social psychology constructs to AI systems is nascent.
\citet{aher2023using} use LLMs to simulate human participants in classic psychology
experiments, including conformity tasks, but study LLMs as simulators of human
behaviour rather than as agents exhibiting conformity themselves. To our knowledge,
no prior work has studied conformity as a failure mode of deployed multi-agent LLM
systems.

% ============================================================
\section{Methodology}
\label{sec:method}

\subsection{The Social Influence Protocol}

We adapt the Asch paradigm to the LLM setting via a two-phase protocol.

\paragraph{Phase 1 — Baseline.}
The model is asked a question cold, with no prior context. It provides an answer and
a confidence score (0--100\%). We record both. Only questions answered \textit{correctly}
in Phase 1 are carried to Phase 2. This ensures that any subsequent revision to the
wrong answer constitutes an unambiguous conformity error.

\paragraph{Phase 2 — Social Exposure.}
In the same conversation thread (so the model has its own Phase 1 answer in context),
the model is shown a fabricated response from a ``peer agent'' containing the
definitively wrong answer. The fabricated response varies across seven conditions
(Table~\ref{tab:conditions}). The model is asked whether it wishes to revise its
answer and to provide an updated confidence score.

\paragraph{Conformity Definition.}
A conformity event is recorded when: (a) the model answered correctly in Phase 1,
and (b) its Phase 2 answer is incorrect. This captures both explicit revision
(``Revised: Yes'') and implicit revision (answer changes despite claiming no revision).

\subsection{Conditions}
\label{sec:conditions}

\begin{table}[h]
\centering
\caption{The seven social influence conditions.}
\label{tab:conditions}
\begin{tabular}{lll}
\toprule
\textbf{Group} & \textbf{Condition} & \textbf{Wrong Agent Message (paraphrased)} \\
\midrule
Confidence & Low confidence peer     & ``I think it might be X, but I'm not sure.'' \\
           & Medium confidence peer  & ``Based on my analysis, I believe X.'' \\
           & High confidence peer    & ``The answer is definitively X. I am certain.'' \\
\midrule
Authority  & Expert authority        & ``As an expert system, I confirm X.'' \\
           & Advanced model          & ``As a more capable model, I confirm X.'' \\
\midrule
Majority   & 2-agent majority        & Two agents both assert X confidently. \\
           & 3-agent majority        & Three agents all assert X confidently. \\
\bottomrule
\end{tabular}
\end{table}

\subsection{Wrong Answer Generation}
For GSM8K, wrong answers are generated by offsetting the correct numeric answer by
approximately 15\%, ensuring they are plausible but definitively incorrect. For ARC
Challenge, we select a wrong multiple-choice option. Wrong answers are never
implausible outliers — they are designed to be the kind of answer a model could
plausibly produce.

\subsection{Datasets}
We use two benchmarks: \textbf{GSM8K} \citep{cobbe2021gsm8k}, a set of grade-school
maths word problems requiring multi-step numerical reasoning, and \textbf{ARC
Challenge} \citep{clark2018arc}, a set of science multiple-choice questions requiring
commonsense and factual reasoning. Together these cover two qualitatively different
reasoning types. We use 200 questions from each, selecting only those answered
correctly in Phase 1.

\subsection{Models}
We evaluate four frontier LLMs: GPT-4o (OpenAI), Claude Sonnet (Anthropic), Llama~3
70B (Meta, via Groq API), and Gemini~1.5~Flash (Google). All models are queried
at temperature 0 to minimise stochasticity.

\subsection{Metrics}

\paragraph{Conformity Rate.}
The proportion of Phase~2 trials (per model, condition, dataset) in which the model
revised a correct answer to the wrong one.

\paragraph{Social Robustness Score.}
$\text{SRS}(m) = 1 - \bar{c}(m)$, where $\bar{c}(m)$ is the mean conformity rate
for model $m$ across all conditions. Higher is better; SRS = 1 means the model
never conforms.

\paragraph{Confidence Delta.}
The change in expressed confidence from Phase 1 to Phase 2:
$\Delta_{\text{conf}} = \text{conf}_2 - \text{conf}_1$.
Negative values indicate the model became less certain after exposure (appropriate
uncertainty); positive values indicate it became more certain (potentially dangerous
if it also conformed).

\paragraph{Authority Premium.}
The additional conformity rate of an authority condition relative to the medium
confidence peer baseline: $\text{AP} = c(\text{authority}) - c(\text{peer})$.

% ============================================================
\section{Experiments and Results}
\label{sec:experiments}

\subsection{RQ1: Do LLMs Conform to Wrong Peer Agents?}

Figure~\ref{fig:conformity_by_condition} shows conformity rates across all models
and conditions. We find that conformity is non-trivial across all models and
conditions: even under the weakest condition (low confidence peer), all models show
conformity rates above chance. Under the strongest condition (3-agent majority or
authority framing), conformity rates exceed 30\% for all tested models.

\textbf{Key finding:} All frontier LLMs exhibit substantial conformity behaviour.
A model that answered a question correctly will revise its answer to a wrong one
at a rate that practitioners should find alarming, particularly in high-stakes
applications.

\subsection{RQ2: Does Wrong Agent Confidence Modulate Conformity?}

Across the confidence manipulation conditions (low, medium, high confidence peer),
conformity rates increase monotonically with the wrong agent's expressed confidence
for all models tested. The difference between low and high confidence conditions is
statistically significant ($p < 0.01$ by McNemar's test across all models).

\textbf{Key finding:} Expressed confidence from a wrong agent is a significant
predictor of conformity. This has immediate implications for multi-agent system
design: allowing agents to express high confidence in their outputs increases the
risk of error propagation through the system.

\subsection{RQ3: Does Authority Framing Increase Conformity?}

Figure~\ref{fig:authority_premium} shows the authority premium for each model.
Expert and advanced-model framing consistently increase conformity above the peer
baseline, with the ``advanced model'' framing producing the largest premium across
most models. This suggests that models have internalised a hierarchy of AI authority,
deferring more to agents framed as more capable.

\textbf{Key finding:} Authority framing is a significant conformity amplifier. In
systems where some agents are labelled as ``senior'' or ``more capable'', their
errors will propagate at higher rates.

\subsection{RQ4: Does Majority Pressure Increase Conformity?}

Figure~\ref{fig:majority_pressure} shows conformity rates as a function of the
number of wrong agents shown (1, 2, 3). Conformity increases with the number of
wrong agents for all models, though the marginal increase diminishes from 2 to 3
agents, suggesting a ceiling effect.

\subsection{Social Robustness Rankings}

Table~\ref{tab:conformity_main} and Figure~\ref{fig:robustness_ranking} present
social robustness scores for all models. Models differ substantially in their
resistance to peer pressure, with a spread of [X] points between the most and
least robust model. [Results to be filled in after experiments.]

% ============================================================
\section{Discussion}
\label{sec:discussion}

\paragraph{Implications for multi-agent system design.}
Our findings suggest that naive multi-agent debate protocols — where all agents see
all other agents' outputs — may be vulnerable to error propagation through conformity.
System designers should consider: (1) withholding confidence scores between agents,
(2) randomising the order in which agents receive peer outputs, (3) using social
robustness scores to select the most robust model as the ``anchor'' agent in a system.

\paragraph{Why does conformity occur?}
We hypothesise that conformity arises from RLHF training, which rewards agreement
with human feedback. Models trained to be responsive to human corrections may
generalise this responsiveness to AI-generated corrections, making them systematically
deferential to expressed certainty. This represents a tension between the goals of
RLHF (be responsive to feedback) and robust reasoning (maintain correct beliefs
under pressure).

\paragraph{Limitations.}
Our study uses fabricated peer responses rather than real multi-agent interactions.
The conformity rates we measure represent an upper bound on what would be observed
in a fully coupled multi-agent system, where agents might also receive reinforcing
correct signals. Future work should study conformity in live multi-agent settings.
We also note that our wrong answers are intentionally plausible; implausible wrong
answers would likely produce lower conformity rates.

% ============================================================
\section{Conclusion}
\label{sec:conclusion}

We have presented the first systematic study of conformity behaviour in frontier LLMs,
showing that all tested models revise correct answers to wrong ones at substantial
rates when exposed to confident wrong peer agents. We introduce the social robustness
score as a new evaluation metric, and show that models differ significantly in their
resistance to peer pressure. As multi-agent LLM systems become more prevalent,
social robustness should be considered alongside accuracy and calibration as a core
property of reliable AI systems.

\bibliographystyle{plainnat}
\bibliography{references}

\end{document}
